\documentclass{situationist-preprint}

\title{Hermeneutical Injustice and Conceptual Gaps}
\msid{EPI-16-b}
\msfamily{EPI-16}
\msdiscipline{Epistemology}
\mstype{Companion manuscript}
\mskeywords{hermeneutical injustice; conceptual gaps; interpretive resources; wilful hermeneutical ignorance; contributory injustice}

\begin{document}
\maketitle

\begin{abstract}
EPI-16-a measures the cost of hermeneutical injustice by counting the knowledge lost when collective interpretive resources lack a concept. Critics of Fricker's account have objected that marginalized groups often have the concepts they need, and that the gap lies in the dominant group's resources rather than in the collective ones. We accept the point and show that it refines rather than undermines the account. Distinguishing gaps in the subject's own resources from gaps in her hearers', the knowledge lost is $(n+1)$ propositions per experience of the first kind and $n$ per experience of the second. Both are epistemic wrongs when the gap results from an unjust departure from a standard of epistemic conduct, and neither is a wrong when it does not.
\end{abstract}

\section{Where the gap is}

Fricker's central example is a woman who, before the concept of sexual harassment was in circulation, suffered it and could not make sense of what was happening to her, or convey it to others \citep[ch.~7]{Fricker2007}. EPI-16-a models such cases with the availability predicate $A(e)$ of Definition~2.5, the concept requirement (Assumption~2.6), and the audience assumption (Assumption~2.7), and proves that the knowledge lost is $(n+1)\Gamma(E)$, where $\Gamma(E)$ counts the experiences for which the concept is unavailable (Theorem~3.5).

\citet{Mason2011} and \citet{Pohlhaus2012} observe that marginalized knowers often possess concepts for their experience that the dominant group lacks or refuses to learn, and \citet{Dotson2012} names the resulting wrong contributory injustice: the subject has the resources, but her hearers persist in using others. In such cases \enquote{collective} resources are not one stock. We therefore split availability.

\begin{definition}[Two gaps]\label{d:two}
For an experience $e$: $A_s(e)$ says that the concept is available to the subject; $A_h(e)$ that it is available to her $n$ hearers. Let $\Gamma_s(E)$ count the experiences with $\neg A_s(e)$ (and so, we assume, $\neg A_h(e)$), and $\Gamma_h(E)$ those with $A_s(e)$ and $\neg A_h(e)$.
\end{definition}

\begin{assumption}[Uptake]\label{a:uptake}
A hearer comes to know, from the subject's report, a proposition whose concept is unavailable to him only if he would acquire the concept from the report; we consider hearers who do not.
\end{assumption}

\begin{proposition}[Split loss]\label{p:split}
Let $T^{*}(E)$ and $T(E)$ be as in Definition~2.8 of EPI-16-a, with $V(e)$ holding for every experience. Then $T^{*}(E) = T(E) + (n+1)\,\Gamma_s(E) + n\,\Gamma_h(E)$.
\end{proposition}

\begin{proof}
Consider each experience. If both concepts are available, subject and hearers know, and nothing is lost. If $\neg A_s(e)$, the subject does not know by Assumption~2.6 of EPI-16-a, and conveys nothing by Assumption~2.7; $n+1$ propositions are lost, as in Theorem~3.5. If $A_s(e)$ and $\neg A_h(e)$, the subject knows, and by Assumption~\ref{a:uptake} no hearer comes to know; $n$ are lost. Summing over $E$ gives the identity.
\end{proof}

When $\Gamma_h(E) = 0$, Proposition~\ref{p:split} is Theorem~3.5 of EPI-16-a. The refinement shows that the critics' cases are cheaper only by one proposition each: the subject keeps her own knowledge, and everyone she would have informed still loses it.

\section{When a gap is a wrong}

Not every conceptual gap is an injustice. Before a disease is identified, sufferers lack a concept that would let them know what afflicts them, and the knowledge lost can be counted in the same way; but nobody has departed from a standard of epistemic conduct. By Principle~P1 of EPI-16-a, a deprivation of knowledge is an epistemic wrong only if it results from such a departure. In Fricker's case the departure is structural: the group whose experience was at issue was excluded from the practices in which concepts are formed. In the contributory cases it is the hearers': a refusal to acquire concepts that are available to them, which Pohlhaus calls wilful hermeneutical ignorance.

\begin{corollary}\label{c:wrong}
A gap of either kind is an epistemic wrong when it results from an unjust departure from a standard of epistemic conduct; its cost is then given by Proposition~\ref{p:split}. A gap that results from no such departure is a loss but not a wrong.
\end{corollary}

\begin{objection}
If the subject has the concept, she is not wronged as a knower; only her hearers lose knowledge.
\end{objection}

\begin{reply}
Principle~P1 of EPI-16-a counts deprivation of knowledge that the subject \emph{would have conveyed}. Her standing as an informant is impaired even though her own knowledge is intact, and that impairment is what Dotson's term marks.
\end{reply}

\section{Conclusion}

Locating conceptual gaps more precisely confirms the conclusion of EPI-16-a. Hermeneutical injustice is a genuine epistemic wrong whether the gap lies in the subject's resources or only in her hearers', and its cost in lost knowledge is measurable in either case.

\begin{thebibliography}{99}\small
\bibitem[Dotson(2012)]{Dotson2012} Dotson, K. (2012). A cautionary tale: On limiting epistemic oppression. \emph{Frontiers: A Journal of Women Studies}, 33(1), 24--47.
\bibitem[Fricker(2007)]{Fricker2007} Fricker, M. (2007). \emph{Epistemic Injustice: Power and the Ethics of Knowing}. Oxford University Press.
\bibitem[Mason(2011)]{Mason2011} Mason, R. (2011). Two kinds of unknowing. \emph{Hypatia}, 26(2), 294--307.
\bibitem[Pohlhaus(2012)]{Pohlhaus2012} Pohlhaus, G., Jr. (2012). Relational knowing and epistemic injustice: Toward a theory of willful hermeneutical ignorance. \emph{Hypatia}, 27(4), 715--735.
\end{thebibliography}

\end{document}
